\section{Introduction}
\label{sec:introduction}

Vector search underpins recommendation~\cite{youtube_dnn}, semantic search~\cite{fb_ebr}, and retrieval-augmented generation (RAG)~\cite{rag}.
As collections grow to billions of vectors, keeping the full index in memory becomes prohibitively expensive.
Disk-based graph indexes such as DiskANN~\cite{diskann} address this by keeping only product-quantized vectors in DRAM and storing each node's full vector and adjacency list on SSD, packing several nodes into a 4\,KB page.
A search walks the proximity graph by repeatedly reading the page of the next node, of which typically only that one node is useful~\cite{starling}---a high read amplification.
For example, on the SIFT1B graph~\cite{bigann}, DiskANN issues 119 disk reads to visit 121 nodes per query on average, and I/O accounts for about 90\% of query latency.
Reducing the number of random disk reads is therefore the most direct way to improve the performance of disk-based graph search.

A disk read pays off only when the fetched page also holds nodes that the query will visit next---in a graph traversal, the neighbors of the node just read.
Proximity graphs give every node dozens of neighbors, while a 4\,KB page holds only a handful of nodes, so no single-copy layout can place a node beside all of its neighbors, and replicating every neighborhood would multiply the index several times over.
Existing systems therefore decide, before serving and from the index alone, which nodes to keep in memory and how to arrange them on disk.
In memory, DiskANN caches nodes reachable by BFS from the entry point, and Laser~\cite{laser} caches high-in-degree nodes.
On SSD, Starling~\cite{starling} and MARGO~\cite{margo} rearrange pages by graph topology so that neighbors tend to share a block.
Run-time buffer pools follow actual accesses, but they cache only individual nodes and are bounded by DRAM.
As a result, neither the memory cache nor the on-disk layout reflects which nodes and edges queries actually traverse, and neither adapts when the workload changes.

Real search workloads, however, need only a small part of this locality.
Queries are often skewed~\cite{quake,ragcache,xie_locality_2002}: many of them concentrate on a few regions of the index, which we call \emph{hotspots}.
Within a hotspot, a small set of \emph{hot nodes} is read repeatedly, and of the dozens of edges leaving each hot node, queries keep following the same few---its \emph{hot edges} (\S\ref{sec:background}).
Caching hot nodes in memory removes their disk reads, but DRAM holds only the hottest of them; the remaining reads still follow those hot edges.
Which edges are hot is decided by where queries go, not by the graph topology or vector distances from which existing caches and layouts are derived.
Consequently, a BFS or degree-based memory cache barely helps under repetitive queries, while topology-driven page rearrangements cannot put a hot node beside the neighbors that searches will actually need.

We therefore use the SSD itself to install denser pages into memory.
We present \textsc{NeighborCache}, which tracks node and edge hotness while serving queries, keeps the hottest nodes in an in-memory \emph{N-cache}, and maintains an SSD-side \emph{D-cache}: each entry is centered on a seed node together with neighbors reached through its hot edges (using as many 4\,KB pages as that payload needs), so a miss brings several soon-needed nodes into memory and later hops hit there.
Every node can be a seed; we cap materialized D-cache space (e.g., 20\% of the index) and leave the index and search algorithm untouched.
As hotspots shift, stale hotness decays and pages are replaced online.
The gain is fewer random reads and higher memory-hit efficiency after each read, so the D-cache can live on the same SSD as the index.

Realizing this idea poses three challenges.
First, maintenance must be cheap: many query threads update node and edge hotness concurrently, and synchronization can erase the I/O savings.
Second, online page replacement must not disturb queries: a reader may be about to touch a page that is being rewritten, so the system must expose only complete pages and must not block other queries on the write.
Third, the caches must follow shifting hotspots at low cost: if hotness only accumulates, past hotspots never yield their pages; old accesses must lose influence without locking all metadata on every decay.

\textsc{NeighborCache} addresses them as follows.
It shards hotness metadata and batches per-query updates at the end of each query, recording edge hotness only from a sampled fraction of queries.
It writes a new D-cache page to a free location and publishes it only after the write completes, reclaiming the old page after readers drain, so queries always see complete pages and writers hold no lock during I/O.
Instead of periodically lowering every counter, it gradually increases the weight of new accesses as the global access count grows---equivalent to decaying old hotness---and rescales counters only when a value approaches the limit of its data type, so D-cache pages of fading hotspots are replaced by nodes in the new ones.

We implement \textsc{NeighborCache} in DiskANN and evaluate it on SIFT100M, SIFT1B, and additional datasets under skewed, repetitive, and hotspot-shifting workloads, comparing against DiskANN with a matched BFS memory budget and against Starling and MARGO.
At the same recall, \textsc{NeighborCache} improves throughput by up to 84\% over DiskANN and reduces disk reads per query by up to 55\%; even on the dataset's nearly uniform query set it already nearly doubles DiskANN, and skewed or repetitive workloads widen the gap further.

We make the following contributions:
\begin{itemize}
  \item We show that selectively caching hot nodes and hot edges already nearly doubles DiskANN on the dataset's nearly uniform query set, that skew amplifies this gain, and that topology- and distance-based caches and layouts miss these edges (\S\ref{sec:background}).
  \item We show that replicating only hot nodes and the neighbors reached through their hot edges, in SSD space bounded to 20\% of the index, overcomes the limit of single-copy layouts without the cost of full neighborhood replication.
  \item We design \textsc{NeighborCache}, which maintains an in-memory N-cache and SSD-resident D-cache pages online from observed searches, with low maintenance cost, safe concurrent page replacement, and decaying hotness so that the cache follows shifting hotspots (\S\ref{sec:design}).
  \item We implement \textsc{NeighborCache} in DiskANN and demonstrate its effectiveness under skewed and hotspot-shifting workloads (\S\ref{sec:eval}).
\end{itemize}
